\section{Q1: History Encodings and Diagnostics}
\label{app:q1}

\subsection{Reference Conditions and Encoding Comparison}

Q1 asks how finite history should be encoded to support motion understanding,
future prediction, and action. A history encoding is evaluated under a declared
future target and world--action reference, rather than under an assumed final
SE(3)/flow pipeline. The phase-specific reference, encoder candidates retained
after screening, input fields, history span, and sampling intervals are
\placeholder{MISSING: Q1 comparison configuration}. The formal result table
must identify whether the compared encoders see the same samples and differ
only in organization, or whether an explicit sampling/window axis also changes.

A current-image control is distinguished from a strict memory-free control.
If the policy retains end-effector history, past actions, a recurrent state, or
previous predicted chunks, it still has memory even when visual history is
removed. The historical label H0V therefore denotes a visual-history
intervention only; its retained robot/action fields must be listed. GT object
history is an information diagnostic and appears in a separate track from
visual-history encodings.

\begin{table}[t]
  \centering
  \caption{Q1 configuration record. Placeholder rows reserve separate
  comparisons; they do not assert that a candidate has been implemented or run.}
  \label{tab:app_q1_configuration}
  \small
  \begin{tabularx}{\columnwidth}{@{}lXX@{}}
    \toprule
    Comparison & Required common conditions & Unresolved configuration \\
    \midrule
    Encoding organization & Same input samples, future target, action objective & \Pending \\
    Window/sampling & Same visible fields, declared physical span and gaps & \Pending \\
    Visual-history removal & List retained EEF/action memory & \Pending \\
    Strict memory-free & No hidden recurrence or cached history & \Pending \\
    GT-history diagnostic & Matched GT track; never pooled with visual SR & \Pending \\
    \bottomrule
  \end{tabularx}
\end{table}

\subsection{Motion Ambiguity, Occlusion, and Sampling}

Exact-current comparisons examine whether histories distinguish motion that is
ambiguous in a current frame. A direction diagnostic matches present pose and
appearance while changing the preceding motion. The construction must be
audited before learning results are interpreted: residual pose, visibility,
or robot-state cues could otherwise explain a difference. Accepted pair counts
and current-state matching tolerances are \Pending.

Occlusion comparisons report the observed history interval, occlusion duration,
and visibility at the query and predicted execution time. Removing older
frames and changing temporal spacing are separate manipulations. The study
reports seconds as well as frame count, and keeps the action horizon fixed
where the question concerns only history. For irregular histories, actual
time offsets and masks are retained instead of relabeling samples at uniform
intervals.

History swapping tests the effect of substituted temporal evidence with current
inputs held fixed. A useful outcome must connect changes in predicted motion,
the selected action, task success, and safety in the direction implied by the
paired scene. Probe predictability or action sensitivity alone cannot show
that historical information caused a useful action improvement. The
history-necessary subset is validated independently of the aggregate SR;
its acceptance evidence is \Pending.

\subsection{Expanded Results and Reporting}

Table~\ref{tab:app_q1_results} reserves the expanded summary. The final table
is indexed by actual encoder, task, and seed rather than by invented candidate
names. For each comparison, the per-task/seed source, number of matched
episodes, interval calculation, and failure exclusions are \Pending.

\begin{table}[t]
  \centering
  \caption{Pending Q1 results. SR differences use the corresponding matched
  reference within a track; GT and visual inputs remain separate.}
  \label{tab:app_q1_results}
  \small
  \begin{tabularx}{\columnwidth}{@{}Xccc@{}}
    \toprule
    Track/diagnostic & SR (\%) & $\Delta$ SR & Unsafe (\%) \\
    \midrule
    Visual encoding comparison & \Pending & \PendingCI & \Pending \\
    Strict memory-free comparison & \Pending & \PendingCI & \Pending \\
    Direction-ambiguous subset & \Pending & \PendingCI & \Pending \\
    Occluded subset & \Pending & \PendingCI & \Pending \\
    Sampling-gap comparison & \Pending & \PendingCI & \Pending \\
    GT-history diagnostic & \Pending & \PendingCI & \Pending \\
    \bottomrule
  \end{tabularx}
\end{table}

Prediction errors are reported at the physical times used for action, with
translation and orientation distinguished when those quantities exist.
Action outcomes additionally include contact-time error, relative pose at
contact, safety, and task completion. Their available result arrays are
\placeholder{MISSING: Q1 per-task, per-seed and per-horizon results}.
No positive effect of longer history or a particular encoder is established
by the present placeholders.

The selected Q1 encoding is provisional until Q2 determines which future
information is useful. After changing the future representation, retained
history candidates are revisited under the new reference; the completed
cross-checks and any ranking changes are reported in Appendix~\ref{app:q3}.
